\begin{center}
    {\LARGE \textbf{2011物理化学}} \\[2ex]
    {\large 时量180分钟 \quad 满分150分}
\end{center}

\vspace{0.5cm}
\noindent\textbf{注意事项:}
\begin{enumerate}[label=\arabic*., parsep=0pt, itemsep=0pt]
    \item 答题前,考生先将自己的姓名、学号、考场号、座位号填写在答题卡上,将条形码准确粘贴在考生信息条形码粘贴区。
    \item 考试结束后,将本试卷与答题纸一并收回。
    \item 回答各个问题时,将答案写在答题纸上,写在本试卷上无效。
\end{enumerate}

\vspace{0.5cm}
\noindent\textbf{一、选择题(30 pts.)}

\begin{enumerate}[label=\arabic*., itemsep=1.5ex]
    \item 在温度为$T$、压力为$p$时,反应$\ce{3O2(g) -> 2O3(g)}$的$K_p$与$K_x$的比值为
    \begin{tasks}(4)
        \task $RT$
        \task $p$
        \task $(RT)^{-1}$
        \task $p^{-1}$
    \end{tasks}
    \item 某化学反应在$298\ \mathrm{K}$时的标准吉布斯自由能变化为正值,则该温度时反应的$K_p^\ominus$将是
    \begin{tasks}(4)
        \task $K_p^\ominus=1$
        \task $K_p^\ominus=0$
        \task $K_p^\ominus>1$
        \task $K_p^\ominus<1$
    \end{tasks}
    \item 已知挥发性纯溶质$\ce{A}$液体的蒸气压为$67\ \mathrm{Pa}$,纯溶剂$\ce{B}$的蒸气压为$26\ 665\ \mathrm{Pa}$,该溶质在此溶剂的饱和溶液的摩尔分数为$0.02$,则此饱和溶液(假设为理想液体混合物)的蒸气压为
    \begin{tasks}(4)
        \task $600\ \mathrm{Pa}$
        \task $26\ 133\ \mathrm{Pa}$
        \task $26\ 198\ \mathrm{Pa}$
        \task $599\ \mathrm{Pa}$
    \end{tasks}
    \item 理想气体在等温条件下反抗恒定外压膨胀,该变化过程中体系的熵变$\Delta S_\text{体}$及环境的熵变$\Delta S_\text{环}$应为
    \begin{tasks}(2)
        \task $\Delta S_\text{体}>0,\Delta S_\text{环}=0$
        \task $\Delta S_\text{体}<0,\Delta S_\text{环}=0$
        \task $\Delta S_\text{体}>0,\Delta S_\text{环}<0$
        \task $\Delta S_\text{体}<0,\Delta S_\text{环}>0$
    \end{tasks}
    \item 含有非挥发性溶质$\ce{B}$的水溶液,在$101\ 325\ \mathrm{Pa}$下,$270.15\ \mathrm{K}$时开始析出冰,已知水的$K_\mathrm{f}=1.86\ \mathrm{K\cdot kg\cdot mol^{-1}}$,$K_\mathrm{b}=0.52\ \mathrm{K\cdot kg\cdot mol^{-1}}$,该溶液的正常沸点是
    \begin{tasks}(4)
        \task $370.84\ \mathrm{K}$
        \task $372.31\ \mathrm{K}$
        \task $373.99\ \mathrm{K}$
        \task $376.99\ \mathrm{K}$
    \end{tasks}
    \item 饱和溶液中溶剂的化学势$\mu$与纯溶剂的化学势$\mu^*$的关系式为
    \begin{tasks}(4)
        \task $\mu=\mu^*$
        \task $\mu>\mu^*$
        \task $\mu<\mu^*$
        \task 不能确定
    \end{tasks}
    \item 下列电池中,电动势与氯离子活度无关的电池是
    \begin{tasks}(1)
        \task $\ce{Zn | ZnCl2(aq) \parallel KCl(aq) | AgCl | Ag}$
        \task $\ce{Pt | H2 | HCl(aq) | Cl2 | Pt}$
        \task $\ce{Ag | AgCl(s) | KCl(aq) | Cl2 | Pt}$
        \task $\ce{Hg | Hg2Cl2(s) | KCl(aq) \parallel AgNO3(aq) | Ag}$
    \end{tasks}
    \item 已知：
    
    $\ce{Tl+ + e^- -> Tl(s)}$,\quad $E_1\ (\ce{Tl+ / Tl})=-0.34\ \mathrm{V}$

    $\ce{Tl^3+ + 3e^- -> Tl(s)}$,\quad $E_2\ (\ce{Tl^3+ / Tl})=0.72\ \mathrm{V}$

    则$\ce{Tl^3+ + 2e^- -> Tl+}$的$E_3$值为
    \begin{tasks}(4)
        \task $1.06\ \mathrm{V}$
        \task $0.38\ \mathrm{V}$
        \task $1.25\ \mathrm{V}$
        \task $0.83\ \mathrm{V}$
    \end{tasks}
    \item 用补偿法(对消法)测定可逆电池的电动势时,主要为了
    \begin{tasks}(1)
        \task 消除电极上的副反应
        \task 减少标准电池的损耗
        \task 在可逆情况下测定电池电动势
        \task 简便易行
    \end{tasks}
    \item $1\ \mathrm{mol}$单原子分子理想气体,从$273\ \mathrm{K}$、$202.65\ \mathrm{kPa}$,经$pT=$常数的可逆途径压缩到$405.3\ \mathrm{kPa}$的终态,该气体的$\Delta U$为
    \begin{tasks}(4)
        \task $1\ 702\ \mathrm{J}$
        \task $-406.8\ \mathrm{J}$
        \task $406.8\ \mathrm{J}$
        \task $-1\ 702\ \mathrm{J}$
    \end{tasks}
    \item 在$298\ \mathrm{K}$时,气相反应$\ce{H2 + I2 -> 2HI}$的$\Delta_\mathrm{r}G_\mathrm{m}^\ominus=-16\ 778\ \mathrm{J\cdot mol^{-1}}$,则反应的平衡常数$K_p^\ominus$为
    \begin{tasks}(4)
        \task $2.0\times10^{12}$
        \task $5.91\times10^{6}$
        \task $873$
        \task $18.9$
    \end{tasks}
    \item $298\ \mathrm{K}$时,$1\ \mathrm{mol}$理想气体的平动能近似为
    \begin{tasks}(4)
        \task $600\ \mathrm{J\cdot mol^{-1}}$
        \task $1\ 250\ \mathrm{J\cdot mol^{-1}}$
        \task $2\ 500\ \mathrm{J\cdot mol^{-1}}$
        \task $3\ 719\ \mathrm{J\cdot mol^{-1}}$
    \end{tasks}
    \item 反应$\ce{2N2O5 -> 4NO2 + O2}$的速率常数单位是$\mathrm{s^{-1}}$。对该反应的下述判断哪个对?
    \begin{tasks}(4)
        \task 单分子反应
        \task 双分子反应
        \task 复合反应
        \task 不能确定
    \end{tasks}
    \item 从热力学基本关系式得知$(\partial A/\partial V)_T$等于
    \begin{tasks}(4)
        \task $(\partial H/\partial S)_p$
        \task $(\partial G/\partial T)_p$
        \task $(\partial H/\partial T)_S$
        \task $(\partial U/\partial V)_S$
    \end{tasks}
    \item 两个活化能不相同的反应,如$E_2>E_1$,且都在相同的升温区间内升温,则
    \begin{tasks}(2)
        \task $\displaystyle \frac{\mathrm{d}\ln k_2}{\mathrm{d}T}>\frac{\mathrm{d}\ln k_1}{\mathrm{d}T}$
        \task $\displaystyle \frac{\mathrm{d}\ln k_2}{\mathrm{d}T}<\frac{\mathrm{d}\ln k_1}{\mathrm{d}T}$
        \task $\displaystyle \frac{\mathrm{d}\ln k_2}{\mathrm{d}T}=\frac{\mathrm{d}\ln k_1}{\mathrm{d}T}$
        \task $\displaystyle \frac{\mathrm{d}k_2}{\mathrm{d}T}>\frac{\mathrm{d}k_1}{\mathrm{d}T}$
    \end{tasks}
\end{enumerate}

\vspace{0.5cm}
\noindent\textbf{二、计算题(15 pts.)}

$1\ \mathrm{mol}$氮气从$100\ \mathrm{K}$,$4.1\ \mathrm{dm^3}$加热到$600\ \mathrm{K}$,$49.2\ \mathrm{dm^3}$,其反抗$100\ \mathrm{kPa}$的恒定外压以不可逆方式进行,且气体可视为理想气体。试计算该体系的$Q$,$W$,$\Delta U$,$\Delta H$,$\Delta S_\text{体}$,$\Delta S_\text{环}$,$\Delta S_\text{孤}$。

\vspace{0.5cm}
\noindent\textbf{三、计算题(20 pts.)}

在$101.325\ \mathrm{kPa}$时,甲苯的沸点为$383.2\ \mathrm{K}$,已知此条件下$\Delta S_\mathrm{m}=87.03\ \mathrm{J\cdot K^{-1}\cdot mol^{-1}}$。试计算:
\begin{enumerate}[label=(\arabic*), itemsep=1ex]
    \item 在上述条件下,$1\ \mathrm{mol}$液态甲苯蒸发时的$Q$,$W$,$\Delta U_\mathrm{m}$,$\Delta H_\mathrm{m}$,$\Delta A_\mathrm{m}$,$\Delta G_\mathrm{m}$;
    \item 欲使甲苯在$373.2\ \mathrm{K}$时沸腾,其相应的饱和蒸气压为多少?
\end{enumerate}

\vspace{0.5cm}
\noindent\textbf{四、计算题(15 pts.)}

在$85^\circ\mathrm{C}$、$101.3\ \mathrm{kPa}$时,甲苯($A$)和苯($B$)组成的液态混合物达沸腾。试计算该理想液态混合物的液相及气相的组成。已知苯的正常沸点为$80.10^\circ\mathrm{C}$,甲苯在$85.00^\circ\mathrm{C}$时的蒸气压为$46.00\ \mathrm{kPa}$。

\vspace{0.5cm}
\noindent\textbf{六、计算题(15 pts.)}

由乙烷裂解制乙烯反应:
\[
\ce{CH3CH3(g) -> CH2=CH2(g) + H2(g)}
\]
在$1\ 000\ \mathrm{K}$和$151.988\ \mathrm{kPa}$下,标准平衡常数$K_p^\ominus=0.898$,反应开始前体系中只有$2\ \mathrm{mol}$乙烷,求:
\begin{enumerate}[label=(\arabic*), itemsep=1ex]
    \item 反应达平衡时的反应进度;
    \item 乙烯的最大产率;
    \item 平衡混合物中各气体的摩尔分数。
\end{enumerate}

\vspace{0.5cm}
\noindent\textbf{七、计算题(15 pts.)}

在$671\sim768\ \mathrm{K}$之间,$\ce{C2H5Cl}$气相分解反应$\ce{(C2H5Cl -> C2H4 + HCl)}$为一级反应,速率常数$k(s^{-1})$和温度($T$)的关系式为：

    {\centering
    $\lg(k/s^{-1})=-13\ 290/(T/\mathrm{K})+14.6$ \par
    }

\begin{enumerate}[label=(\arabic*), itemsep=1ex]
    \item 求$E_\mathrm{a}$和$A$;
    \item 在$700\ \mathrm{K}$时,将$\ce{C2H5Cl}$通入一反应器中($\ce{C2H5Cl}$的起始压力为$26\ 664.5\ \mathrm{Pa}$),反应开始后,反应器中压力增大,问需多少时间反应器中压力变为$46\ 662.8\ \mathrm{Pa}$?
\end{enumerate}

\vspace{0.5cm}
\noindent\textbf{八、计算题(20 pts.)}

$298\mathrm{K}$时有电池:
$\ce{Pt,H2(p) | KOH(a=1),Ag2O(s) | Ag(s)}$,

已知:

$E(\ce{Ag2O | Ag})=0.344\ \mathrm{V}$,

$\Delta_\mathrm{f}H_\mathrm{m}^\ominus(\ce{H2O,l})=-286.0\ \mathrm{kJ\cdot mol^{-1}}$, $\Delta_\mathrm{f}H_\mathrm{m}^\ominus(\ce{Ag2O,s})=-30.57\ \mathrm{kJ\cdot mol^{-1}}$。

试求:
\begin{enumerate}[label=(\arabic*), itemsep=1ex]
    \item 电池的电动势,已知$K_\mathrm{w}=1.0\times10^{-14}$;
    \item 电池反应的$\Delta_\mathrm{r}G_\mathrm{m}^\ominus$, $\Delta_\mathrm{r}H_\mathrm{m}^\ominus$, $\Delta_\mathrm{r}S_\mathrm{m}^\ominus$, $Q_\mathrm{r}$ 和$\left(\dfrac{\partial E}{\partial T}\right)_p$;
    \item 若电池在短路下放电,程度与可逆反应相同时,则热效应为多少?
\end{enumerate}